In \cite{synthetic-stone-duality}, we show that all functions on the unit interval are continuous in the $\epsilon-\delta$ sense. In this section, we will extend these results to uniform continuity. We will do so in two ways. One way is to show that continuity implies uniform continuity. The other way is to use the notion of uniform space, as in \parencite[Ch2]{BourbakiTopologyCh1To4}. 



The following is a proof using the topology on $\mathbb I$. 
\begin{lemma}[uniform continuity of maps $\mathbb I \to \mathbb R^n$]\label{uniformContinuity}
  For any map $f:\mathbb I \to [0,1] \times [-1,1]$, and for every $\epsilon:\mathbb R$ with $\epsilon>0$, 
  there exists a $\delta>0$ such that for any $x,y:\mathbb I$ with $d(x,y)<\delta$, we have
  that $d(f(x) , f(y))<\epsilon$.
\end{lemma}
\begin{proof}
  Because $f$ is continuous, for each $x:\mathbb I$ there exists some $n:\mathbb N$ such that 
  if $d(x,y)<\frac1n$, then $d((f(x), f(y)))<\epsilon$. 
  By local choice, we get a Stone space $S$ with a surjection $a:S\twoheadrightarrow \mathbb I$ 
  and a map $b:S\to \mathbb N$ such that if $d(x,a(s))<\frac1{b(s)}$, then $d(f(x), f(a(s)))<\epsilon$. 
  By Lemma 3.6 of \cite{synthetic-stone-duality}, $b$ factors over $Fin(k)$. 
  Hence we can take the maximal element $N$ of it's image, and set $\delta = \frac1N$. 
\end{proof}
\begin{remark}
  By LLPO and dependent choice, the Cauchy reals agree with the reals in the setting of SSD. Thus we can define any function which comes with a Cauchy sequence. In particular, in the setting of SSD, we can define the following functions:
  \begin{itemize}
    \item For each $k:\mathbb N$, we can define $t_k: [\frac1k , 1] \to [-1,1]$ by $t_k(r) = \sin(\frac1r)$. 
    \item We can define for each $n:\mathbb N$ the standard metric $d:\mathbb R^n \times \mathbb R^n \to \mathbb R_{\geq 0}$. 
  \end{itemize}
\end{remark}






